\chapter{Introduction}

Corporate bankruptcy prediction has emerged as a critical area of research in financial analytics, offering substantial value to stakeholders including investors, creditors, and policymakers. The ability to accurately predict bankruptcy enables informed decision-making, risk assessment, and proactive financial management strategies. This study leverages a comprehensive bankruptcy dataset spanning multiple time horizons to develop robust ensemble models for bankruptcy prediction.

The dataset utilized in this research comprises bankruptcy information for companies across five different time periods (from 1 to 5-year horizons). Each temporal dataset contains 64 financial features that capture various aspects of company performance, including liquidity ratios, profitability metrics, leverage indicators, and operational efficiency measures. The complete dataset encompasses thousands of companies, providing a rich foundation for developing predictive models.

The financial features in the dataset represent key financial ratios and indicators commonly used in financial analysis and bankruptcy prediction. These include metrics such as working capital ratios, debt-to-equity ratios, return on assets, current ratios, and various profitability measures. The comprehensive nature of these features allows for a holistic assessment of company financial health across multiple dimensions.

The primary objective of this research is to develop and compare ensemble machine learning models for bankruptcy prediction using the consolidated multi-year dataset. Specifically, this study aims to:

\begin{enumerate}
    \item Merge the five temporal datasets (1-year through 5-year horizons) into a unified dataset for comprehensive analysis
    \item Conduct thorough exploratory data analysis to understand data characteristics, missing value patterns, and feature relationships
    \item Implement data preprocessing techniques including missing value imputation and feature selection strategies
    \item Develop and optimize two ensemble models: Random Forest and XGBoost classifiers
    \item Evaluate model performance using both full feature sets and selected feature subsets
    \item Compare model effectiveness through comprehensive evaluation metrics including ROC curves, AUC scores, and classification reports
    \item Provide insights into the most predictive features and optimal modeling approaches for bankruptcy prediction
\end{enumerate}

The choice of ensemble methods, specifically Random Forest and XGBoost, is motivated by their proven effectiveness in handling complex financial data, their ability to capture non-linear relationships, and their robustness to overfitting. These tree-based ensemble methods are particularly well-suited for bankruptcy prediction tasks due to their capacity to handle mixed data types, identify important features, and provide interpretable results.

This research contributes to the growing body of literature on bankruptcy prediction by providing a comprehensive comparison of ensemble methods across different feature selection strategies. The findings will offer practical insights for financial analysts and practitioners seeking to implement effective bankruptcy prediction systems in real-world applications.

The subsequent chapters present detailed exploratory data analysis, methodology and modeling approaches, and comprehensive results with discussions on model performance and practical implications for bankruptcy prediction in financial decision-making contexts. 